%==============================================================================
% Antonio Balanzategui — Resume
% Rebuilt as an editable LaTeX source (recreated from Antonio_Balanzategui_Resume_1.pdf).
% Compile with:  tectonic Antonio_Balanzategui_Resume.tex
%==============================================================================
\documentclass[letterpaper,9pt]{extarticle}

\usepackage[T1]{fontenc}
\usepackage[utf8]{inputenc}
\usepackage{lmodern}
\usepackage[margin=0.35in]{geometry}
\addtolength{\textheight}{0.30in}
\linespread{0.93}
\usepackage{enumitem}
\usepackage{titlesec}
\usepackage[hidelinks]{hyperref}
\usepackage{xcolor}

\pagestyle{empty}
\setlength{\parindent}{0pt}

% --- Section headings: bold, small caps, rule underneath -----------------------
\titleformat{\section}
  {\normalsize\bfseries\scshape}
  {}{0em}{}[\vspace{-0.6ex}\rule{\linewidth}{0.6pt}]
\titlespacing{\section}{0pt}{1.0ex}{0.5ex}

% --- Bullet list tuned tight ---------------------------------------------------
\setlist[itemize]{leftmargin=1.4em, topsep=0.2ex, itemsep=0.05ex, parsep=0pt,
  label=\textbullet}

% --- Helper: an entry with a right-aligned date --------------------------------
\newcommand{\entry}[2]{%
  \textbf{#1}\hfill#2\par}
\newcommand{\subentry}[2]{%
  \textit{#1}\hfill\textit{#2}\par\vspace{0.2ex}}

\begin{document}

%==============================================================================
% HEADER
%==============================================================================
\begin{center}
  {\LARGE\bfseries Antonio Balanzategui}\\[0.6ex]
  Richmond, VA 23235 $\vert$ (804) 332-9563 $\vert$ \href{mailto:antbalanzategui@vt.edu}{antbalanzategui@vt.edu}\\[0.3ex]
  \href{https://linkedin.com/in/antbalanzategui}{linkedin.com/in/antbalanzategui} $\vert$
  \href{https://github.com/antbalanzategui}{github.com/antbalanzategui} $\vert$
  \href{https://antoniobalanzategui.vercel.app}{antoniobalanzategui.vercel.app}
\end{center}

%==============================================================================
% EXPERIENCE
%==============================================================================
\section{Experience}

\entry{Project Engineer --- Control Systems Software}{Jan 2026 -- Present}
\subentry{Engineered Services, Inc. (contracted to VCU Health)}{Richmond, VA}
\begin{itemize}
  \item Initial February statistical analysis correctly predicted a mechanical change-point from BAS sensor data alone (CUSUM, Mann-Kendall); validated three months later when the EF1 exhaust fan failed catastrophically on March 19, 2026. VCU Health Pharmacy regulatory compliance leadership formally requested quarterly investigations going forward.
  \item Write and revise production sequence-of-operations logic on the hospital's BAS (EcoStruxure Building Operation, Andover Continuum) for RTUs, exhaust fans, pumps, pharmacy cleanrooms, and air-compressor failover.
  \item Programmed a new burn-ICU expansion's rooftop unit and eight VAV boxes from scratch through startup and test-and-balance: an AS-P state-machine program (lockouts, changeover, freeze/smoke shutdown, dual-fan proof, static-pressure reset, min-OA flow) and native MP-V VAV control with a safety guard. Proved every branch in 14 drills on the live controller; fixed a PID anti-windup defect and wrote a custom XPI block that sets each PID's mode from unit state and control mode so idle loops cannot wind up.
  \item Diagnose and retune unstable control loops with frequency-domain analysis: modeled a chilled-water pump loop as first-order-plus-dead-time, found the $-180^\circ$ phase crossover, and showed the installed gain left a gain margin below 1 --- a one-line gain change stopped pumps hunting 25--100\% and held $\Delta$P at setpoint in production. In separate self-directed analysis, proved by reachability argument and an MPC benchmark that no controller could hold a fan-coil unit's commanded state, then designed a damped positional PI replacement ($\zeta$ 0.13$\to$0.55, \textasciitilde13$\times$ less valve travel in simulation). Both written up as public case studies.
  \item Built and maintain an automated Python pipeline over 10+ months of BAS trend-log data across 12 sensor channels --- change-point detection, multivariate SPC, Hawkes processes, Kalman smoothing, Weibull survival --- generating the full quarterly investigation end-to-end and replacing manual analysis.
  \item Diagnosed a BAS historian stale-data cascade (Hawkes branching ratio n* = 0.973, near-critical) across all 12 channels, tracing it to an ExTL buffer / Maximum Transfer Interval mismatch; the config fix cut worst-case staleness 24 hours $\to$ 1.
  \item Built two internal tools on the MaintainX and Freshdesk REST APIs: one that migrates service tickets out of MaintainX into Freshdesk, replacing manual re-keying of hundreds of tickets; and one that auto-fills service tickets from Freshdesk time-log entries and routes them to sales. Both in use across the campus.
\end{itemize}

\vspace{0.6ex}
\entry{Undergraduate Teaching Assistant, Introduction to GUI Programming}{Aug 2024 -- May 2025}
\subentry{Virginia Tech, Department of Computer Science}{Blacksburg, VA}
\begin{itemize}
  \item Taught 100+ students object-oriented design, event-driven programming, GUI architecture, and the mathematics of 2D graphics.
\end{itemize}

%==============================================================================
% PROJECTS
%==============================================================================
\section{Projects}

\entry{EvolutionSim --- Real-Time GPU Thunderstorm \& Lightning Engine}{2025 -- Present}
\subentry{Personal research project (private repo); 40-page methods \& verification paper and demo available on request}{}
\begin{itemize}
  \item Built a real-time GPU storm engine (Rust + wgpu + WGSL, \textasciitilde100 storm compute kernels) simulating deep moist convection coupled to explicit electrification on grids up to 5.3M cells: conservative flux-form (MPDATA) advection, virtual-potential-temperature buoyancy, geometric-multigrid pressure \emph{and} electric-field solvers, two-moment bulk microphysics, and non-inductive graupel--ice charge separation.
  \item Made the GPU hybrid the production engine for the paper's storm results, with the double-precision CPU implementation kept as a kernel-parity oracle (per-kernel gates at \textasciitilde$10^{-6}$ relative RMS on the gated build, i.e. single-precision tolerance); profiled the flash handler --- the host round-trip I meant to optimize was 1.6\% of per-flash time, the multigrid re-solve 67\% --- then parallelized that solve's red-black smoother across 16 cores for a 1.9$\times$ \emph{bit-identical} speedup, rejecting a faster variant that perturbed the result. 5--8$\times$ the single-threaded reference at matched resolution; four real-sounding storms complete in 6--18 min each in one unattended overnight queue.
  \item Made lightning emergent rather than scripted: field-directed stepped leaders grow along the resolved electric field, and each flash is typed cloud-to-ground / intracloud / cloud-air by where its channel terminates --- the charge tripole, lightning jump, and multicell self-organization all arise from the physics; flash rates land within 1.38--1.86$\times$ of Price--Rind.
  \item Checked against observation: total-water conservation to machine precision on the f64 reference ($10^{-9}$\,relative on the GPU); 14 real radiosonde soundings (CAPE 0--4822 J/kg) reproducing the continental-vs-maritime updraft contrast; and two severe-storm soundings --- matching the Norman, OK tornadic storm's NEXRAD echo top to \textasciitilde4\% (15.3 vs \textasciitilde15.9 km), and a Dodge City storm whose depth, intensity and electrification are consistent with the large hail reported in the same air mass; an automated 9-check capability scorecard scores every run from saved output alone, and a six-seed ensemble separates the exactly-deterministic fluid/microphysics solve from the one stochastic component (mean flash rate robust to \textasciitilde4\%). Part of a 371-mode platform with an \textasciitilde10{,}000-test pytest suite.
\end{itemize}

%==============================================================================
% EDUCATION
%==============================================================================
\section{Education}

\entry{Bachelor of Science, Computer Science}{May 2025}
\subentry{Virginia Tech, Blacksburg, VA}{}
\begin{itemize}
  \item GPA: 3.3 $\vert$ Pollard Scholarship $\vert$ 2nd Place, VTURCS Undergraduate Research Symposium --- de novo protein-structure prediction from sequence via simulated annealing in AMBER (Scientific Computing).
\end{itemize}

%==============================================================================
% SKILLS
%==============================================================================
\section{Skills}

\textbf{Languages:} Python, Rust, C++, Java, JavaScript, TypeScript, SQL, C, R, MATLAB; Spanish (proficient).\par\vspace{0.25ex}
\textbf{GPU \& Graphics:} wgpu, WGSL, compute shaders, GPU-accelerated PDE solvers, CPU$\leftrightarrow$GPU parity validation, volumetric rendering.\par\vspace{0.25ex}
\textbf{Scientific Computing \& Physics:} Navier-Stokes (Boussinesq, Chorin projection), conservative flux-form advection (MPDATA), geometric multigrid PDE solvers, stencil methods on structured grids, two-moment bulk microphysics, sub-grid turbulence closures (Deardorff), CRM verification (Kessler, Weisman--Klemp), molecular dynamics (AMBER).\par\vspace{0.25ex}
\textbf{ML, Statistics \& Data:} NumPy, SciPy, pandas, scikit-learn, PyTorch, LightGBM, matplotlib, data pipelines; CUSUM, PELT, Mann-Kendall, MEWMA Hotelling T-squared, Hawkes processes, Kalman/RTS smoothing, Weibull survival, PCA.\par\vspace{0.25ex}
\textbf{Software \& Systems:} pytest, parity harnesses \& golden-state regression, profiling, REST APIs, Docker, Linux, Git, PostgreSQL, LaTeX.\par\vspace{0.25ex}
\textbf{Control \& BAS:} PID/PI tuning, frequency-domain stability (gain \& phase margin), dead-time compensation, anti-windup, MPC; EcoStruxure, Andover Continuum, BACnet/IP \& MSTP.

\end{document}
